%=============================================================================
\section{Conclusion}
%=============================================================================

Le système constitue une plateforme fonctionnelle pour le calcul automatisé de potentiels redox par DFT. Les principales caractéristiques sont : (i)~un pipeline conformères en 4 étapes avec GFN1-xTB et DFT+freq+solvant obligatoire, (ii)~un schéma composite TZ2P+DZP pour le pipeline redox assurant la cohérence des bases et un résiduel de décomposition nul, (iii)~la réutilisation du neutre optimisé au lieu de le recalculer, et (iv)~un module d'analyse vibronique comprenant 6 méthodes. Les pipelines sont opérationnels et ont été appliqués à 73 molécules dans 10 solvants, produisant 1411 conformères exploitables après filtrage.

La décomposition thermodynamique de $\Delta G$ en quatre termes (EA/IP, $E_\text{relax}$, $\Delta\Delta U$, $T\Delta S$), désormais exacte (résiduel $= 0$) grâce au schéma composite, révèle une hiérarchie claire. En oxydation, IP et $E_\text{relax}$ représentent plus de 99\% de $|\Delta G|$. En réduction, EA et $E_\text{relax}$ restent dominants mais les corrections thermiques sont plus significatives.

Le modèle prédictif multi-solvant combine l'approximation de Koopmans (IP, EA), le couplage vibronique linéaire LVC2 ($E_\text{relax}$), et des descripteurs CDFT ($\Delta\Delta U$, $T\Delta S$) avec des coefficients $\varepsilon$-dépendants (approximants de Padé). Le modèle prédit $\Delta G_\text{ox}$ avec MAE = \SI{6.0}{kJ/mol} ($R^2_\text{adj} = 0{,}967$) et $\Delta G_\text{red}$ avec MAE = \SI{5.7}{kJ/mol} ($R^2_\text{adj} = 0{,}994$), sans optimisation géométrique ionique. L'eau ($\varepsilon = 78{,}4$) reste le cas limite (MAE $\sim$ \SI{18}{kJ/mol}) en raison des limites du modèle COSMO implicite.

Les prochaines étapes prioritaires sont : (i)~la réduction du résidu anharmonique sur $E_\text{relax}$ pour les molécules à forte réorganisation conformationnelle, (ii)~l'extension à d'autres familles redox-actives, et (iii)~la validation expérimentale lors du stage à l'Université de York.

